\uuid{pQWk}
\titre{Raisonnements}
\auteur{bodin ; croizat ; sacre}
\organisation{exo7}
\licence{CC BY-NC-SA 4.0}
\datecreate{2022-03-08}
\niveau{L1}
\module{Algèbre}
\chapitre{Logique, ensemble, raisonnement}
\sousChapitre{Absurde et contraposée}
\difficulte{1}
\theme{Raisonnements}
\exo7qcmid{hilisit.logique.044}

\contenu{
\begin{question}
\qclassification{100.03}
On commence une démonstration par l'absurde avec la rédaction suivante : "Supposons que $\log_{10}(3) \in \Qq$. Alors on peut écrire $\log_{10}(3) = \frac p q$ avec ...". Que cherche-t-on à démontrer ?
\begin{answers}
	\bad{$\log_{10}(3) \in \Qq$}
	\good{$\log_{10}(3) \notin \Qq$}
	\bad{$\log_{10}(3) \in \Rr$}
	\bad{$\log_{10}(3) \notin \Rr$}
\end{answers}
\begin{explanations}
On a commencé par supposer que $\log_{10}(3) \in \Qq$. Puisqu'il s'agit d'une démonstration par l'absurde, cette supposition correspond à la négation de ce qu'on doit démontrer. On doit donc démontrer que $\log_{10}(3) \notin \Qq$.
\end{explanations}
\end{question}
}
